\documentclass{beamer}

\usepackage{amsmath}
\usepackage{amssymb}
\usepackage{booktabs}
\usepackage{graphicx}
\usepackage{xcolor}
\usepackage{makecell}

\title{Model comparison and diagnosis}
\author{Jiahui Fan}
\date{\today}

\graphicspath{{.}{pred_vs_raw_compare/}}

\begin{document}

\frame{\titlepage}

% ----------------------------------------------------------------
\begin{frame}{Data setup: main simulation}
\small
\textbf{Simulated data dimensions:}
\begin{tabular}{ll}
\toprule
Dimension & Value \\
\midrule
Participants ($P$)                & 20 \\
EEG electrodes ($C$)              & 61 \\
Conditions ($K$) (noise; noise + signal) & 2 \\
Time points ($T$)                 & 51 (0--500\,ms, 100\,Hz) \\
Trials per participant per condition ($R$) & 20 \\
\bottomrule
\end{tabular}

\bigskip
\textbf{Raw data:} $P \cdot C \cdot K \cdot T \cdot R = 2{,}488{,}800$ rows.

\smallskip
\textbf{Aggregated data} (grouped by ppt, electrode, condition, time):
$P \cdot C \cdot K \cdot T = 124{,}440$ rows.
\end{frame}

% ----------------------------------------------------------------
\begin{frame}{Ground truth: noise model}
\small
Each trial has additive noise:
\[
\text{noise} \;=\; \text{temporal noise} \;+\; \text{spatial noise}
\]

\medskip
\textbf{Temporal 1/f noise} (Yeung et al., 2004): 50 sinusoids
with EEG-like power spectrum,
$\mathrm{Var} = \sigma^{2}_{\text{out}} = 0.5$, drawn independently
for each (trial, electrode).
\end{frame}

% ----------------------------------------------------------------
\begin{frame}{Ground truth: signal model}
\small
\[
\text{signal}
\;=\;
\underbrace{\text{peak amplitude}}_{A\,=\,1.5\,\mu\mathrm{V}}
\;\cdot\;
\underbrace{\text{temporal shape}}_{g(z)}
\;\cdot\;
\underbrace{\text{spatial weight}}_{w(c)}
\]

\medskip
Truncated Gaussian $g(z)$, with $z = (t - \tau_p)/\sigma_t$ and
$\sigma_t = 0.060$\,s, $\tau_p \approx 0.250$\,s:
\[
g(z) = \frac{\phi(z) - \phi(1.5)}{\phi(0) - \phi(1.5)},
\quad |z| \le 1.5; \quad g(z) = 0 \text{ else.}
\]
$\phi$ = standard Gaussian density. $g(0)=1$, $g(\pm 1.5)=0$;
window $\approx [0.160, 0.342]$\,s.

\medskip
Spatial Gaussian RBF (unnormalised Gaussian, peak $=1$ at CZ):
\[
w(c) = \exp\!\Bigl(-\,\frac{d^{2}(c, \mathrm{CZ})}{2\,\sigma_{\mathrm{space}}^{2}}\Bigr),
\quad \sigma_{\mathrm{space}} = 0.5.
\]

\end{frame}

% ----------------------------------------------------------------
\begin{frame}{Spatial ground truth: active electrodes}
\small
The current simulation uses a hard-threshold spatial Gaussian.

\[
d(c,\mathrm{CZ}) = \left\|p_c - p_{\mathrm{CZ}}\right\|_2
\]

\[
w(c) =
\begin{cases}
\exp\!\left(-\dfrac{d^2(c,\mathrm{CZ})}
{2\sigma_{\mathrm{space}}^2}\right),
& d(c,\mathrm{CZ}) \leq \sigma_{\mathrm{space}},\\[0.7em]
0,
& d(c,\mathrm{CZ}) > \sigma_{\mathrm{space}}.
\end{cases}
\]

\medskip
\[
\text{electrode } c \text{ is active iff }
d(c,\mathrm{CZ}) \leq \sigma_{\mathrm{space}}.
\]

\medskip
This is the ground-truth active-electrode list used for Spatial F1,
Spatio-temporal F1, FP, and FN metrics.
\end{frame}

% ----------------------------------------------------------------
\begin{frame}{Simulated EEG data ($P=20$, $R=20$)}
\begin{center}
\includegraphics[width=\linewidth]{sim_data_overview_n20_r20.pdf}
\end{center}
\end{frame}

% ----------------------------------------------------------------
\begin{frame}[fragile]{Modeling}
\scriptsize
\begin{verbatim}
bgam_fit <- brm(
    formula = eeg_mean | se(eeg_sd) ~
        condition +
        (1 + condition | participant) +
        t2(xproj, yproj, time,
           by   = condition,
           bs   = c("tp", "tp"),
           k    = c(20, 20),
           d    = c(2, 1),
           full = FALSE),
    data    = sim_data_model,
    family  = gaussian(),
    chains  = 4,
    cores   = 4,
    iter    = 3000,
    warmup  = 1000,
    seed    = 42,
    threads = threading(threads = 4),
    backend = "cmdstanr",
    control = list(adapt_delta = 0.99, max_treedepth = 13)
)
\end{verbatim}
\end{frame}

% ----------------------------------------------------------------
\begin{frame}{How the model fits the data ($P=20$, $k=(20,20)$)}
\begin{center}
\includegraphics[height=0.85\textheight,keepaspectratio]{pred_vs_raw_bgam_fit_v5_k20_n20_r20_sd.pdf}
\end{center}
\end{frame}


% ----------------------------------------------------------------
\begin{frame}{Posterior probability per electrode ($P=20$, $k=(20,20)$)}
\begin{center}
\includegraphics[width=0.92\linewidth]{posterior_p_sig_facets_n20_r20.pdf}
\end{center}
\end{frame}

% ----------------------------------------------------------------
\begin{frame}{Model result: $P=20$, $R=20$, $k=(20,20)$}
\scriptsize
\textbf{Fit:} \texttt{bgam\_fit\_v5\_k20\_n20\_r20\_sd.rds}

\medskip
\begin{center}
\renewcommand{\arraystretch}{1.25}
\begin{tabular}{lcc}
\toprule
Diagnostic target & Result & Interpretation \\
\midrule
Convergence & PASS & 0 divergences; $\widehat{R}_{\max}=1.003$ \\
Onset / offset & 150--350\,ms & truth 160--342\,ms; bias $-10/+8$\,ms \\
Spatial top channel & CZ & correct seed electrode \\
Spatial F1 & 0.847 & 25 TP, 9 FP, 0 FN \\
Spatiotemporal F1 & 0.831 & 343 TP, 8 FP, 132 FN \\
Peak recovery & 76.9\% & pointwise peak is underestimated \\
Window-mean recovery & 105.6\% & 95\% CrI covers truth \\
Marginal condition diff. & 0.0529 & empirical diff. $=0.0533$ \\
Wall clock & 173\,h & total execution time on LIDO HPC \\
\bottomrule
\end{tabular}
\end{center}

\medskip
This result covers the temporal window, spatial localisation, and
window-averaged amplitude. The pointwise peak is still underestimated.
\end{frame}

% ----------------------------------------------------------------
\begin{frame}{Topomaps over the signal window ($P=20$, $R=20$)}
\begin{center}
\includegraphics[height=0.86\textheight,keepaspectratio]{results/figures/topomap_n20_r20_three_times.pdf}
\end{center}
\end{frame}

% ----------------------------------------------------------------
\begin{frame}{SE calibration: effective trial count}
\scriptsize
\textbf{Data/model:} $P=5$ participants, $R=20$ trials per condition,
$61$ electrodes, $51$ time points, $k=(20,20)$.

\medskip
\textbf{Definitions for one participant $\times$ channel $\times$ condition $\times$ time cell:}
\[
\mathrm{SD} = \text{trial-level variability},
\qquad
\mathrm{SEM} = \frac{\mathrm{SD}}{\sqrt{R}},
\qquad
\mathrm{SE}_{\mathrm{eff}} = \frac{\mathrm{SD}}{\sqrt{R_{\mathrm{eff}}}}.
\]
\vspace{-0.5em}
\[
R_{\mathrm{eff}} \leq R
\quad\text{is the effective number of independent trials after trial dependence.}
\]

\smallskip
\textbf{Here:} $R_{\mathrm{eff}}=14$ comes from an empirical trial-correlation
ESS check in the active window/active channels ($13.77 \approx 14$).
$R_{\mathrm{eff}}=10$ is a more conservative intermediate sensitivity value.

\medskip
\begin{center}
\renewcommand{\arraystretch}{1.18}
\resizebox{\linewidth}{!}{%
\begin{tabular}{lccccccccc}
\toprule
Fit & $R_{\mathrm{eff}}$ & \makecell{Onset\\bias} & \makecell{Offset\\bias} &
\makecell{Peak\\rec.} & \makecell{Mean\\rec.} & \makecell{Spatial\\F1} &
\makecell{ST\\F1} & LOOic & \makecell{Wall\\clock} \\
\midrule
\texttt{k20\_reff10} & 10 & 0\,ms & $-12$\,ms & 86.4\% & 105.7\% & 0.806 & 0.814 & 1289.4 & 69\,h \\
\texttt{k20\_reff14} & 14 & 0\,ms & $-12$\,ms & 88.1\% & 105.9\% & 0.794 & 0.816 & $-2311.5$ & 68\,h \\
\bottomrule
\end{tabular}}
\end{center}

\medskip
Both fits recover the temporal window and CZ-centred spatial effect. Peak is
still underestimated, but window-mean amplitude is close to truth.
\end{frame}

% ----------------------------------------------------------------
\begin{frame}{Methods comparison ($P=5$ pilot, 30 MC)}
\scriptsize
Following Nalborczyk \& Bürkner (2025) Table 1.
All rows are 30-MC summaries from the $P=5$ pilot setting.
Onset/offset in seconds; SD across 30 sims.

\medskip
\textbf{Onset estimates} (truth $= 0.160$\,s):

\begin{tabular}{lrrrrc}
\toprule
Method        & Bias (s)            & MAE (s)             & RMSE (s)            & SD (s)             & $F1_{\text{ST}}$ (mean\,$\pm$\,SD) \\
\midrule
Raw $p$       & $+0.010$            & $\mathbf{0.010}$    & $0.012$             & $0.008$            & $0.462 \pm 0.014$ \\
\textbf{BGAM\_2D} & $-0.011$        & $0.011$             & $\mathbf{0.012}$    & $\mathbf{0.005}$   & $\mathbf{0.799 \pm 0.007}$ \\
TFCE          & $+0.013$            & $0.014$             & $0.017$             & $0.011$            & $0.539 \pm 0.105$ \\
Cluster       & $+0.016$            & $0.016$             & $0.017$             & $0.007$            & $0.514 \pm 0.020$ \\
FDR (BH)      & $+0.063$            & $0.063$             & $0.065$             & $0.016$            & n/a \\
Holm          & $+0.093$            & $0.093$             & $0.094$             & $0.012$            & n/a \\
\bottomrule
\end{tabular}

\medskip
\textbf{Offset estimates} (truth $= 0.342$\,s):

\begin{tabular}{lrrrr}
\toprule
Method        & Bias (s)            & MAE (s)             & RMSE (s)            & SD (s) \\
\midrule
\textbf{BGAM\_2D} & $+0.012$        & $\mathbf{0.012}$    & $\mathbf{0.014}$    & $\mathbf{0.006}$ \\
Raw $p$       & $-0.006$            & $0.013$             & $0.015$             & $0.014$ \\
TFCE          & $-0.015$            & $0.015$             & $0.018$             & $0.011$ \\
Cluster       & $-0.018$            & $0.019$             & $0.020$             & $0.009$ \\
FDR (BH)      & $-0.068$            & $0.068$             & $0.070$             & $0.017$ \\
Holm          & $-0.082$            & $0.082$             & $0.082$             & $0.010$ \\
\bottomrule
\end{tabular}

\end{frame}

% ----------------------------------------------------------------
\begin{frame}{FP and FN definitions}
\scriptsize
This checks the previous onset/offset summary metrics using the full
time-by-channel detection mask.

\smallskip
Each cell is one time-by-channel location $(t,c)$.

The ground-truth mask is:
\[
M_{t,c} =
\begin{cases}
1, & t \in [0.160, 0.342]\text{ s and }c\text{ is an active electrode},\\
0, & \text{otherwise}.
\end{cases}
\]

\smallskip
A detected cell is a cell declared significant by a given method.
For detection, we use $p_{t,c}<\alpha$, where $\alpha = 0.10$.

\smallskip
Let $D_{t,c}=1$ if cell $(t,c)$ is detected, and $D_{t,c}=0$ otherwise.

\[
\begin{array}{ll}
TP:\ D_{t,c}=1,\ M_{t,c}=1, & FP:\ D_{t,c}=1,\ M_{t,c}=0,\\
FN:\ D_{t,c}=0,\ M_{t,c}=1, & TN:\ D_{t,c}=0,\ M_{t,c}=0.
\end{array}
\]

\[
\text{FP rate} = \frac{FP}{FP + TN},
\qquad
\text{FN rate} = \frac{FN}{FN + TP}.
\]
\end{frame}

% ----------------------------------------------------------------
\begin{frame}{FP/FN diagnosis ($\alpha=0.10$)}
\scriptsize
Onset/offset metrics only summarise the first and last
detected time points. They do not show whether the full detection mask
contains many false-positive cells or misses most true cells.

\medskip
\begin{tabular}{lrrrrr}
\toprule
Method & FP rate & FN rate & $F1_{\text{ST}}$ & \makecell{Mean FP\\cells} & \makecell{FP-only\\clusters} \\
\midrule
Raw $p$  & $0.102$ & $0.530$ & $0.462$ & $268.9$ & $141.4$ \\
BH       & $0.001$ & $0.943$ & $0.106$ & $3.0$   & $2.9$ \\
Holm     & $0.000$ & $0.996$ & $0.008$ & $0.2$   & $0.2$ \\
Cluster  & $0.003$ & $0.648$ & $0.514$ & $8.0$   & $0.0$ \\
TFCE     & $0.007$ & $0.608$ & $0.539$ & $18.4$  & $1.1$ \\
\bottomrule
\end{tabular}

\medskip
Raw $p$ has the highest FP rate because each time-channel cell is tested
separately without correction, creating many scattered detections
($141.4$ FP-only clusters; $113.0$ outside-window clusters).
BH and Holm strongly control false positives, but become too conservative
and miss most true cells. Cluster and TFCE use spatio-temporal structure,
so they reduce scattered FP cells and give a better FP/FN trade-off.

\end{frame}

% ----------------------------------------------------------------
\appendix

\begin{frame}{Appendix: Metric definitions}
\small
$F1$ measures detection accuracy of active cells.

\medskip
\[
F1 \;=\; \frac{2\,P\,R}{P + R},
\quad
P = \frac{TP}{TP + FP},
\quad
R = \frac{TP}{TP + FN}.
\]
$P$ = fraction of predicted-positive cells that are truly active. \\
$R$ = fraction of truly active cells that are predicted positive.

\medskip
A cell is \textbf{predicted positive} iff its posterior odds $> \theta$.
This is the threshold the BGAM uses to declare a cell ``significant'',
and it defines the $TP$, $FP$, $FN$ counts above.

\medskip
The two metrics differ only in what counts as a ``cell'':

\begin{center}
\renewcommand{\arraystretch}{1.3}
\begin{tabular}{lll}
\toprule
                          & \textbf{Spatial F1}             & \textbf{Spatio-temporal F1} \\
\midrule
Cell                      & channel $c$                     & (channel, time) pair $(c, t)$ \\
Evaluated at              & peak time $t^{\star} = 0.251$\,s & all $51$ time points \\
Total cells               & $61$ channels                   & $61 \times 51 = 3{,}111$ \\
True positives            & $25$ active channels             & $25 \times 19 = 475$ active cells \\
\bottomrule
\end{tabular}
\end{center}

\end{frame}

% ----------------------------------------------------------------
\begin{frame}{Appendix: BGAM\_2D odds sensitivity}
\scriptsize
\textbf{Setting:} $P=5$ pilot, 30 Monte Carlo simulations, $v5\_k20$ specification.

\medskip
\begin{center}
\includegraphics[height=0.74\textheight,keepaspectratio]{odds_sensitivity.pdf}
\end{center}

\end{frame}

\begin{frame}{Appendix: BGAM\_2D posterior-mask diagnosis}
\scriptsize
\textbf{Setting:} $P=5$ pilot, 30 Monte Carlo simulations.

\medskip
We define the following rule to obtain the BGAM\_2D posterior mask:
\[
\text{posterior odds}_{t,c}
= \frac{P(\Delta_{t,c}>0 \mid y)}{P(\Delta_{t,c}\le 0 \mid y)}
> 10.
\]
\vspace{-0.4em}
\[
\Delta_{t,c} = \text{condition difference at cell }(t,c)
\quad (\mathrm{cond2}-\mathrm{cond1}).
\]

\medskip
\begin{center}
\resizebox{\linewidth}{!}{%
\begin{tabular}{lrrrrrrrr}
\toprule
Method & \makecell{Odds\\thr.} & \makecell{Mean\\TP} & \makecell{Mean\\FP} &
\makecell{Mean\\FN} & FP rate & FN rate & Precision & $F1_{\text{ST}}$ \\
\midrule
BGAM\_2D (se=SD) & $10$ & $390.1$ & $111.3$ & $84.9$ &
$0.042$ & $0.179$ & $0.778$ & $0.799$ \\
\bottomrule
\end{tabular}}
\end{center}

\medskip
BGAM\_2D has a higher FP rate than cluster and TFCE, but a much lower
FN rate. This means it recovers most of the true spatio-temporal effect,
while also extending somewhat beyond the strict ground-truth mask.

\medskip
This pattern is expected from the smooth Bayesian surface: information
is shared across neighbouring times and channels, which reduces missed
true cells, but can create smooth spillover around the true active region.
\end{frame}

% ----------------------------------------------------------------
\begin{frame}{Appendix: $\alpha$ sensitivity (frequentist methods)}
\begin{center}
\includegraphics[width=0.97\linewidth]{alpha_sensitivity.pdf}
\end{center}

\end{frame}

% ----------------------------------------------------------------
\begin{frame}{Appendix: Larger $P$ improves data-to-parameter ratio}
\footnotesize
\textbf{Spline coefficients} (v5\_k20):
\quad
$\underbrace{20}_{\text{tp loc.}} \times \underbrace{20}_{\text{tp time}} \times \underbrace{2}_{\text{cond.}} = 800.$

\smallskip
\textbf{Signal-bearing observations:}
\quad
$\underbrace{20}_{\text{ppt}} \times \underbrace{25}_{\text{active}} \times \underbrace{19}_{\substack{\text{signal}\\\text{steps}}} \times \underbrace{1}_{\text{cond2}} = 9{,}500.$

\smallskip
\textbf{Ratio:} $9{,}500 / 800 \approx 12$, now within the lower end of the
$\geq 10$--$20$ obs/param rule of thumb (Peduzzi 1996; Wood 2017).

\medskip
With $P=20$, the likelihood carries more information for the same
spatio-temporal smooth. This helps recover the temporal window and spatial
extent, while the smoothness penalty can still attenuate sharp pointwise peaks.

\smallskip
\textit{Remaining caveat:} peak recovery is still affected by smoothing bias.
\end{frame}

% ----------------------------------------------------------------
\begin{frame}{Appendix: Wand--Jones smoothing bias}
\small
Any spline averages each point with its
neighbours over a bandwidth $h$. Near a peak the neighbours
are lower than the peak, so the smoothed estimate sits
\emph{below} the truth. Plugging in our setting ($h \approx 0.025$\,s,
$A = 1.5\,\mu$V, $\sigma_t = 0.060$\,s) gives a peak under-estimate
of $\sim 0.1$--$0.3\,\mu$V $\Rightarrow$ a \textbf{peak-recovery
ceiling of $\sim 85$--$95\%$}, even for an ideal BGAM.

\medskip
\textbf{Formal statement} (Wand \& Jones 1995, Ch.\,4):
\[
E\bigl[\hat f(x_0)\bigr] - f(x_0)
\;\approx\;
\tfrac{1}{2}\, h^2 \, f''(x_0) \, \mu_2(K).
\]
$h$ = effective bandwidth (here $\approx T/k_t = 0.5/20 = 0.025$\,s);
at the peak $f'' < 0$ $\Rightarrow$ negative bias.

\medskip
This hits \emph{Peak recovery} (ceiling $\sim 85$--$95\%$) but is
mostly averaged out in \emph{Mean over window} (ceiling $\sim 95$--$98\%$),
because $f''$ flips sign between centre and edges.
\end{frame}


% ----------------------------------------------------------------
\begin{frame}{Appendix: pilot simulated EEG data ($P=5$, se(eeg\_sd))}
\begin{center}
\includegraphics[width=\linewidth]{sim_data_overview.pdf}
\end{center}
\end{frame}

% ----------------------------------------------------------------
\begin{frame}{Appendix: pilot predictions vs raw ($P=5$, se(eeg\_sd))}
\begin{center}
\includegraphics[height=0.85\textheight,keepaspectratio]{pred_vs_raw_bgam_fit_v5_k20.pdf}
\end{center}
\end{frame}

% ----------------------------------------------------------------
\begin{frame}{Appendix: pilot posterior probability ($P=5$)}
\begin{center}
\includegraphics[width=0.92\linewidth]{posterior_p_sig_facets.pdf}
\end{center}
\end{frame}

% ----------------------------------------------------------------
\begin{frame}{Appendix: pilot model comparison ($P=5$)}
\scriptsize
Original pilot sweep over basis dimensions and model structure.
Here, \textbf{pilot} means the early smaller setting with
$P=5$ participants and $R=20$ trials per condition.

\medskip
\resizebox{\linewidth}{!}{%
\begin{tabular}{ccccrrrrrrr}
\toprule
$k_{xy}$ & $k_t$ & \texttt{full} & SE & \makecell{Onset\\bias (s)} & \makecell{Offset\\bias (s)} & \makecell{Peak\\rec.\,(\%)} & \makecell{Mean\\rec.\,(\%)} & \makecell{Spatial\\F1} & LOOic & \makecell{Wall-\\clock} \\
\midrule
12 & 12 & F                  & SD  & $-0.070$ & $+0.068$ & 46          & 55.4          & $0.714$          & 57462.5          & 5\,h   \\
20 & 20 & F                  & SD  & $-0.070$ & $+0.068$ & 46          & 54.9          & $0.714$          & 57465.5          & 14\,h  \\
\textbf{20} & \textbf{20} & \textbf{F$^{\star}$} & SD & $\mathbf{-0.080}$ & $\mathbf{+0.078}$ & 40 & \textbf{75.5} & $\mathbf{0.758}$ & \textbf{57462.4} & 28\,h \\
\textbf{20} & \textbf{20} & \textbf{F} & \textbf{SEM} & $\mathbf{+0.010}$ & $\mathbf{-0.012}$ & \textbf{89} & \textbf{106.1} & \textbf{0.962} & -- & 67\,h \\
10 & 30 & T                  & SD  & $-0.080$ & $+0.068$ & 36          & 54.3          & $0.714$          & 57469.4          & 48\,h  \\
20 & 30 & T                  & SD  & $-0.130$ & $+0.108$ & 40          & 63.2          & $0.658$          & 57468.7          & 112\,h \\
30 & 20 & T                  & SD  & $-0.150$ & $+0.138$ & 58          & 57.7          & $0.633$          & 57490.4          & 69\,h  \\
30 & 30 & T                  & SD  & $-0.150$ & $+0.138$ & 36          & 57.7          & $0.625$          & 57489.9          & 110\,h \\
\bottomrule
\end{tabular}}

\smallskip
\textsuperscript{$\star$} Winner among SD fits in the original pilot comparison.
\end{frame}

% ----------------------------------------------------------------
\begin{frame}{Appendix: spatial ground-truth alternative}
\small
A less arbitrary alternative is to define the active region by the
half-maximum level set, following the FWHM convention for Gaussian peaks
(\textsuperscript{1}):
\[
\text{electrode } c \text{ is active if } w(c) \geq 0.5\,w_{\max},
\qquad
w_{\max} = \max_{\text{all electrodes } e} w(e).
\]

\medskip
This keeps electrodes within the half-maximum spatial support and is invariant
to changes in the maximum signal amplitude.

\medskip
For EEG electrodes, the final spatial ground truth can combine two parts:
\begin{itemize}\small
\item \textbf{Half-maximum amplitude rule}: define the core active region from
the simulated Gaussian surface.
\item \textbf{Electrode-layout correction}: adjust the grid-based region using
prior knowledge about the electrode layout and expected spatial spread.
\end{itemize}

\vfill
\tiny \textsuperscript{1} Weisstein, Eric W. ``Gaussian Function.''
From \textit{MathWorld}--A Wolfram Resource.
https://mathworld.wolfram.com/GaussianFunction.html. The FWHM section derives
the half-maximum points of a Gaussian and defines the width between them.
\end{frame}

% ----------------------------------------------------------------
\begin{frame}[fragile]{Appendix: Onset / offset bias}
\small
At seed electrode CZ, find times where condition difference is
``significantly positive'' (posterior odds $> 10$, Jeffreys),
then take the temporal extent.

\medskip
\scriptsize
\begin{verbatim}
epred  = add_epred_draws(grid_at_CZ, fit, re_formula = NA)
diff   = cond2 - cond1                       # per draw, per time
p_sig  = mean(diff > 0)                      # per time
odds   = p_sig / (1 - p_sig)

onset  = min{ t : odds(t) > 10 }
offset = max{ t : odds(t) > 10 }

bias_onset  = onset  - 0.160 s
bias_offset = offset - 0.342 s
\end{verbatim}

\smallskip
\textbf{Truth:} $(0.160, 0.342)$\,s (boundary of truncated Gaussian
$g(z)$ where $|z| = 1.5$).
\end{frame}

% ----------------------------------------------------------------
\begin{frame}[fragile]{Appendix: Peak \& Mean recovery (at seed CZ)}
\scriptsize

\textbf{Peak recovery (Target 4) --- point estimate at peak time.}
\begin{verbatim}
peak_t   = (onset + offset) / 2 = 0.251 s
diff_t   = posterior(cond2 - cond1) at (CZ, peak_t)
est_peak = mean(diff_t);  CI = quantile(diff_t, c(.025, .975))
rec_pct  = 100 * est_peak / 1.5     # truth = 1.5 uV (g(0)=1)
\end{verbatim}

\medskip
\textbf{Mean recovery (Target 4b) --- window-mean.}
\begin{verbatim}
window   = [0.160, 0.342] s
# per-draw mean across the window (then aggregate across draws)
mean_d   = per_draw: mean_t( diff(t in window) )
est_mean = mean(mean_d);  CI = quantile(mean_d, c(.025, .975))
rec_pct  = 100 * est_mean / 0.73    # truth = window-mean of g(z)
\end{verbatim}
\end{frame}

% ----------------------------------------------------------------
\begin{frame}[fragile]{Appendix: LOOic}
\small
Bayesian out-of-sample predictive accuracy via PSIS-LOO
(Vehtari, Gelman, Gabry 2017).

\medskip
\scriptsize
\begin{verbatim}
lo  = loo(fit, moment_match = FALSE)

# lo$pointwise[i, "elpd_loo"]  for i = 1 .. 31,110 cells
# elpd_loo_i = log p(y_i | y_{-i})  estimated by PSIS

elpd_loo = sum(lo$pointwise[, "elpd_loo"])
looic    = -2 * elpd_loo
\end{verbatim}

\end{frame}

% ----------------------------------------------------------------
\begin{frame}{Appendix: Topomap (ground truth vs empirical vs model; $P=5$, se(eeg\_sd))}
\begin{center}
\includegraphics[height=0.85\textheight,keepaspectratio]{topoplot_4rows_gt_emp_vs_model.pdf}
\end{center}
\end{frame}

% ----------------------------------------------------------------
\begin{frame}{Appendix: Posterior probability ($k=(20,20)$, \texttt{full=F}; $P=5$, se(eeg\_sd))}
\begin{center}
\includegraphics[width=0.92\linewidth]{posterior_ch_time_heatmap.pdf}
\end{center}
\end{frame}

% ----------------------------------------------------------------
\begin{frame}{Appendix: Posterior log$_{10}$ odds per electrode ($P=5$, se(eeg\_sd))}
\begin{center}
\includegraphics[width=0.92\linewidth]{posterior_log10_odds_facets.pdf}
\end{center}
\end{frame}

% ----------------------------------------------------------------
\end{document}
